\section{Representación lineal, superposición y autoencoder disperso}

\subsection{Features y circuitos: la unidad básica de comprensión}

Muchas neuronas en Inception V1 tienen un significado interpretable claro: curvas, autos, ruedas, orejas de perro. Chris Olah mostró un circuito completo de detección de autos: conectado a un detector de ventanas (arriba), un detector de ruedas (abajo) y un detector de carrocería (en medio) —esto \textbf{es} el algoritmo que detecta autos, leído directamente de los pesos.

Pero el problema es que no todas las neuronas son interpretables. Las \textbf{neuronas polisémicas} (polysemantic neurons) responden a varias cosas no relacionadas entre sí. Esto produce una complejidad exponencial:

\begin{warningbox}{La polisemia produce una explosión exponencial}
Si dos neuronas polisémicas responden cada una a 3 conceptos, el peso entre ellas ya tiene $3 \times 3 = 9$ interacciones posibles que considerar. El problema más profundo es que un espacio de alta dimensión tiene un volumen exponencial: si no puede descomponerse en partes independientes, la complejidad de la comprensión se vuelve inmanejable.

Las \textbf{features monosémicas} (monosemantic features) —features con un único significado claro— permiten razonar de forma independiente y así evitar la explosión exponencial. Por eso perseguir la monosemia es tan importante.
\end{warningbox}

\subsection{La hipótesis de la representación lineal}

\begin{importantbox}{Las direcciones tienen significado}
Las \textbf{direcciones} en el espacio de activaciones tienen significado semántico. Más activación equivale a mayor confianza de detección.

El ejemplo clásico de Word2Vec: King - Man + Woman = Queen; Sushi - Japan + Italy = Pizza. Poder sumar y restar vectores y obtener resultados con significado refleja una estructura lineal.

\textit{«Esto es en realidad lo fundamental que está ocurriendo: las direcciones tienen significado.»} Todo lo observado hasta ahora en redes neuronales naturales es consistente con la hipótesis de la representación lineal.
\end{importantbox}

Chris también mencionó trabajo nuevo sobre representaciones no lineales (features multidimensionales o variedades), pero su postura frente a la hipótesis lineal es pragmática: citó una analogía magnífica:

\begin{knowledgebox}{La lección de la teoría del calórico}
Incluso una teoría equivocada (como la teoría del calórico), si se toma en serio y se lleva hasta el extremo, puede producir resultados prácticos: \textbf{el motor de combustión fue desarrollado por gente que creía en la teoría del calórico}.

\textit{«Vale la pena tomarse en serio una hipótesis y llevarla al límite.»} Aunque la hipótesis de la representación lineal termine por no ser del todo correcta, avanzar dentro de su rango de validez también puede producir descubrimientos importantes.
\end{knowledgebox}

\subsection{La hipótesis de la superposición}

Un embedding de palabra de 500 dimensiones no puede corresponder únicamente a 500 conceptos: el modelo necesita representar muchos más conceptos que dimensiones tiene. La solución: aprovechar las propiedades geométricas del espacio de alta dimensión junto con la dispersión (sparsity) de los conceptos («Japón» e «Italia» rara vez aparecen a la vez).

\begin{importantbox}{«La sombra de una red dispersa más grande»}
\textit{«Las redes neuronales podrían ser la sombra de redes neuronales más grandes y dispersas; lo que vemos son esas proyecciones.»}

Imagina un «modelo del piso de arriba» (upstairs model): enorme pero disperso, donde cada neurona es interpretable. La red neuronal real es una proyección comprimida de ese modelo del piso de arriba. Aprender equivale a construir una compresión eficiente del modelo del piso de arriba. El descenso de gradiente podría estar buscando en secreto el espacio de modelos extremadamente dispersos para luego colapsarlos en una matriz densa.

Cantidad de conceptos: el lema de Johnson-Lindenstrauss muestra que el número de direcciones aproximadamente ortogonales que se pueden embeber es \textbf{exponencial} respecto al número de dimensiones. La dispersión y la estructura de correlación aumentan aún más esta cifra.
\end{importantbox}

\subsection{El avance del autoencoder disperso}

Si la hipótesis de la superposición es correcta, el aprendizaje de diccionario (dictionary learning) mediante autoencoder disperso (sparse autoencoder) es la solución natural.

\textbf{Octubre de 2023}: en «Towards Monosemanticity», al usar un autoencoder disperso en un modelo de una sola capa, \textit{«hermosas features interpretables emergieron así, de forma natural»}. Se encontraron: features de árabe, features de hebreo, features de Base64, features de lenguajes de programación, un «the» específico de cierto contexto (aparece en contextos matemáticos y predice «vector» y «matrix»), y un patrón de alternancia de medios caracteres Unicode.

Al entrenar el modelo dos veces, se encontraron features similares en ambos casos: la universalidad (universality) quedó confirmada de nuevo. \textit{«Una técnica muy natural simplemente funcionó… en realidad es una situación muy buena.»}

\textbf{Mayo de 2024}: se extendió a Claude 3 Sonnet (un modelo de producción). Requirió \textbf{muchísimas} GPU; las leyes de escalamiento para interpretabilidad (Scaling Laws for Interpretability) de Tom Henighan ayudaron a predecir el tamaño óptimo del autoencoder disperso y la cantidad óptima de tokens de entrenamiento.

\begin{warningbox}{Features multimodales abstractas «escalofriantes»}
En Claude 3 Sonnet se encontraron features fascinantes pero también inquietantes:

\begin{itemize}
    \item \textbf{Feature de vulnerabilidad de seguridad}: responde al mismo tiempo al texto «disable SSL» \textbf{y} a imágenes de alguien haciendo clic en una advertencia SSL de Chrome
    \item \textbf{Feature de puerta trasera}: responde al mismo tiempo a puertas traseras (backdoors) en código \textbf{y} a imágenes de cámaras ocultas
    \item \textbf{Feature de engaño o mentira}: al forzar su activación, Claude empieza a mentir
    \item Features de \textbf{búsqueda de poder, golpes de Estado y ocultamiento de información}, entre otras
\end{itemize}

\textit{«Esto muestra qué tan abstractos son estos conceptos.��} La existencia de estas features se relaciona directamente con la necesidad de seguridad de ASL-4 que Dario discutió en la Parte I: se necesita interpretabilidad (interpretability) para detectar si un modelo está siendo engañoso.
\end{warningbox}

\subsection{Los límites de la interpretabilidad automatizada}

Una idea natural es usar a Claude para etiquetar las features que descubre el autoencoder disperso. Chris reconoce el valor de la interpretabilidad automatizada, pero \textit{«tiene ciertas dudas al respecto»}:

Las etiquetas que da la IA \textit{«son ciertas en algún sentido, pero no capturan realmente la feature concreta»}. Esto se parece al escepticismo de los matemáticos frente a las demostraciones automatizadas por computadora: uno cree que la conclusión es correcta, pero no la \textbf{entiende}.

Una preocupación de seguridad más profunda: el artículo clásico de Ken Thompson «Reflections on Trusting Trust» —si usas IA para verificar la seguridad de otra IA, ¿puedes confiar en el auditor? Por ahora no es un problema grande, pero a medida que los sistemas se vuelvan más poderosos, esta pregunta de «quién vigila a los vigilantes» se volverá crucial.

\subsection{Direcciones futuras: de lo microscópico a lo macroscópico}

\textbf{De las features a los circuitos}: hoy se entiende la representación (qué se activa), pero todavía no se entiende el proceso de cómputo (cómo fluye y se transforma la información). Los pesos de interferencia (interference weights, un artefacto de la superposición) hacen más difícil el análisis de circuitos.

\textbf{El problema de la «materia oscura»}: los autoencoders dispersos actuales solo ven una pequeña parte de la «materia» de la red neuronal. \textit{«Es como la astronomía temprana: a medida que construimos mejores autoencoders dispersos… vemos cada vez más estrellas.»}

\textbf{De la microbiología a la anatomía}: la interpretabilidad mecanicista (Mech Interp) actual es «la microbiología de las redes neuronales»: extremadamente fina. Pero las preguntas que nos interesan son macroscópicas. Hace falta ascender de nivel: biología molecular $\to$ biología celular $\to$ histología $\to$ anatomía $\to$ zoología $\to$ ecología. O, como analogía en física: partículas individuales $\to$ física estadística $\to$ termodinámica.

\begin{knowledgebox}{¿Tienen «órganos» las redes neuronales?}
\textit{«Espero que exista algo mucho más grande que las features y los circuitos.»} ¿Existen estructuras macroscópicas análogas al corazón, a regiones cerebrales o al sistema respiratorio?

No se puede saltar directo a lo macroscópico: primero hay que entender la estructura microscópica y luego estudiar los patrones de conexión. Así como no se puede saltar la biología molecular para estudiar directamente la anatomía.
\end{knowledgebox}

\subsection{Redes neuronales artificiales frente a biológicas}

Chris considera que el trabajo de los neurocientíficos es «mucho más difícil». La lista de ventajas de estudiar redes neuronales artificiales es impresionante:

\begin{itemize}
    \item Registrar \textbf{todas} las neuronas (no solo las que se pueden alcanzar)
    \item Introducir datos arbitrarios
    \item Las neuronas no cambian durante el estudio
    \item Se puede ablacionar (eliminar) cualquier neurona
    \item Se puede editar cualquier peso de conexión
    \item Se pueden deshacer todos los cambios
    \item Se puede forzar la activación de cualquier neurona a un valor arbitrario
    \item Se conoce el conectoma completo, \textbf{incluyendo los pesos exactos} (no solo conexiones binarias)
    \item Se pueden calcular gradientes
\end{itemize}

\textit{«Tenemos tantas ventajas frente a los neurocientíficos, y aun con todas estas ventajas, esto sigue siendo muy difícil. Si es tan difícil para nosotros, parece casi imposible bajo las restricciones de la neurociencia.»}

Por eso Chris también recluta activamente a gente del campo de la neurociencia: \textit{«un problema más fácil, pero que sigue siendo muy difícil»}.

\subsection{Seguridad y belleza: el doble llamado de Mech Interp}

Algunas personas se sienten «decepcionadas» de las redes neuronales: \textit{«Solo son la amplificación de reglas simples.»} La réplica de Chris es extraordinariamente elegante: la evolución también «es solo reglas simples»: mutación aleatoria más selección natural. Y aun así produjo todo el esplendor de la biología.

\textit{«La belleza está en que la simplicidad produce complejidad.»} Dentro de las redes neuronales se crea \textit{«una enorme complejidad y belleza que la gente normalmente no ve. Una estructura increíblemente rica que espera ser descubierta.»}

La imagen final: circuitos que crecen hacia la luz de la función de pérdida (loss function) —\textit{«este organismo que hemos cultivado, sin saber qué es lo que hemos cultivado.»} Esto es a la vez el motor científico de Mech Interp y su motor estético.

Cierra con las palabras de Alan Watts: \textit{«The only way to make sense out of change is to plunge into it, move with it, and join the dance.»}

\subsection{Resumen de la sección}

La hipótesis de la representación lineal y la hipótesis de la superposición constituyen la base teórica actual de Mech Interp. El autoencoder disperso extrae features interpretables a partir de la superposición: el éxito al escalar desde un modelo de una sola capa hasta Claude 3 Sonnet demuestra la viabilidad del método. La interpretabilidad automatizada es útil, pero no puede sustituir la comprensión humana. El reto central a futuro es pasar de la «microbiología» a la «anatomía»: descubrir la estructura macroscópica de las redes neuronales. Mech Interp no solo sirve a la seguridad, sino que también explora la pregunta científica profunda de qué es exactamente lo que hemos creado.
